% 對應英文：sections/02_introduction.tex

加密貨幣市場已由 \citet{nakamoto2008bitcoin} 所提出的單一去中心化資產，
發展為一個涵蓋眾多代幣、鏈上經濟機制與政策曝險的異質資產市場。
\citet{liu2022cryptocurrency} 指出，加密貨幣的橫截面報酬具有自身的風險因子結構，
其中動量與規模溢酬雖與股票市場相似，卻在形成機制與市場微結構上呈現明顯差異。
在此背景下，預測哪些代幣會在下一週相對於其他資產表現更佳，
本質上並非單純的時間序列點預測問題，而是一個需要比較資產彼此相對強弱的橫截面排序問題。
對於一個涵蓋 86 種積極交易代幣、同時包含價格、鏈上、總體、政策訊號與 ETF 資金流資料的宇宙而言，
模型不僅需要整合異質特徵，還必須學會一種「相對報酬」而非「絕對報酬」的預測概念。

\paragraph{既有文獻與研究缺口。}
\citet{gu2020empirical} 是機器學習資產定價文獻中的代表性研究，
指出在特徵集充足且模型容量受到適當正則化時，
樹模型與神經網路能夠系統性超越線性基準。
\citet{freyberger2020dissecting} 進一步說明，
資產特徵之間的非線性交互作用本身就攜帶顯著的額外預測力。
這些發現自然引導出一個問題：
若橫截面特徵在股票市場已具有高度解釋力，
那麼在波動更大、資訊來源更雜訊化的加密貨幣市場中，
機器學習是否同樣能有效地從橫截面訊號中擷取可交易的排序資訊？

加密貨幣文獻已經證明此市場具備可預測性，但對模型機制的結論並不一致。
\citet{cakici2024crypto_ml} 發現，
加密貨幣橫截面報酬的可預測性確實存在，
但模型複雜度所帶來的邊際收益可能有限；
\citet{filippou2024crypto_ml} 則顯示，
在涵蓋動量、網路活動、技術指標與投資人關注的豐富特徵集合下，
非線性模型可產生顯著的樣本外預測與經濟價值。
這些研究共同說明，問題不在於加密貨幣是否可預測，
而在於如何讓模型以正確的方式學習可比較、可交易的橫截面信號。

本文認為，深度學習模型在此類任務中經常表現不如預期，
關鍵原因並不一定是架構不足，而是目標函數與評估標準之間的錯配。
當標準 ListNet 損失 \citep{cao2007learning} 直接以原始下週報酬作為 softmax 目標時，
不同週次的目標分布會因波動幅度不同而產生截然不同的尖銳度。
在高波動週次，少數極端資產主導梯度；
在平靜週次，目標分布又接近均勻，導致梯度訊號過弱。
\citet{poh2021building} 在股票橫截面研究中已指出這個問題，
並建議以排名百分位取代原始報酬作為 ListNet 的監督訊號。
本文將此想法延伸至加密貨幣市場，
因其跨資產波動率差異更大，故目標尺度不一致的問題更加嚴重。

\paragraph{資料面向的研究動機。}
除了損失函數外，本文的另一個動機來自資料來源的擴充。
\citet{cong2021tokenomics} 與 \citet{liu2022cryptocurrency} 指出，
鏈上指標如 NVT、MVRV 與網路活動資料，
可反映加密貨幣與股票市場截然不同的基本面結構。
另一方面，傳統金融文獻早已證明替代資料的重要性：
\citet{tetlock2007giving} 顯示媒體語調對股票短期價格有預測力，
\citet{bollen2011twitter} 則發現 Twitter 情緒可領先部分市場變動。
在加密貨幣市場中，\citet{maitre2025social} 進一步指出，
社群媒體注意力與橫截面報酬之間存在可檢測的關聯。
因此，若研究只停留在價格與鏈上資料，可能低估與政策訊號或敘事熱度相關的報酬來源。

本文特別引入兩類過去較少在同一框架下整合的結構性資料。
第一類是川普在 Truth Social 與 X 上的發文特徵，
用以捕捉高影響力政策行為者在關稅與加密監管議題上的公開表態；
第二類是 Farside Investors 自 2024 年 1 月現貨比特幣 ETF 核准後所提供的細粒度資金流資料
\citep{bitcoin_etf_farside}，
可將 ETF 總流量拆解為 Grayscale 存量資金贖回與排除 GBTC/ETHE 後的新增資金流入。
在我們所掌握的文獻中，
將鏈上基本面、政策相關社群媒體訊號與 ETF 結構性資金流
同時放入橫截面排序模型中進行比較，仍屬相對少見。

\paragraph{架構設計上的考量。}
在模型架構方面，本文同時檢驗三種不同的歸納偏誤。
LSTM \citep{hochreiter1997long,fischer2018lstm_finance}
擅長從有限回看窗口中擷取時序依賴，是金融時間序列中穩健的基準模型。
TFT \citep{lim2021temporal} 則透過變數選擇網路與注意力機制，
強調可解釋性與對異質輸入的處理能力。
然而，\citet{gu2020empirical} 的結論也提醒我們：
若真正有效的是橫截面當期特徵，那麼複雜的時序編碼器未必必要。
基於這一點，本文另外提出不含時序編碼器的 CrossSectionalGatedNet，
藉由門控殘差網路與跨資產注意力，直接檢驗「純橫截面架構是否已足以完成排序任務」。

\paragraph{研究貢獻。}
基於上述研究缺口，本文有四項主要貢獻：
\begin{enumerate}
  \item \textbf{排名正規化的 ListNet 目標。}
        依循 \citet{poh2021building}，
        以橫截面排名百分位 $\in [0, 1]$ 取代原始報酬作為 softmax 目標，
        降低不同市場環境下的目標尺度不一致問題。
  \item \textbf{梯度累積。}
        透過將有效橫截面批次大小由每步 86 個資產提升至
        $86 \times 4 = 344$ 個資產，
        進一步穩定排序損失的梯度方向。
  \item \textbf{橫截面門控網路（CrossSectionalGatedNet）。}
        本文提出一個不含時序編碼器的架構，
        由多層 GRN 與跨資產注意力構成，
        用以檢驗橫截面結構本身是否足以支撐加密貨幣的排序任務。
  \item \textbf{修正後的 TFT 變數選擇網路。}
        我們以每變數共享的 GRN 取代過度簡化的輕量 VSN，
        使 TFT 的實作更貼近 \citet{lim2021temporal} 的原始規格，
        並提升特徵選擇的穩定性。
\end{enumerate}

\paragraph{章節架構。}
第~\ref{sec:related}~節回顧相關文獻，
並說明本文與既有加密貨幣機器學習研究的差異。
第~\ref{sec:data}~節介紹資料面板、特徵來源與偏誤控制。
第~\ref{sec:method}~節形式化排序損失與模型架構。
第~\ref{sec:experiments}~節說明超參數設定與六組資訊集設計。
第~\ref{sec:results}~節報告跨模型比較結果。
第~\ref{sec:analysis}~節進一步討論消融研究、Granger 因果檢定、子期間穩健性與集成效果。
第~\ref{sec:conclusion}~節總結研究發現與限制。
